%! Justifying a Claim About the Difference of Two Means Based on a CI — Unit 7, Lesson 7.7 — Guided Notes
\documentclass[10pt]{article}
\usepackage{apstats-article}
\usepackage{apstats-boxes}
\newcommand{\TallMath}[1]{$\displaystyle #1\rule[-1.4em]{0pt}{3.2em}$}

\begin{document}

\pageheader{Unit 7, Lesson 7.7}{Guided Notes}
\namedateperiod

\vspace{0.15in}

\begin{objectivebox}
By the end of this lesson, I will be able to:
\begin{itemize}
  \item \writeline
  \item \writeline
  \item \writeline
\end{itemize}
\end{objectivebox}

\begin{vocabbox}
\begin{description}
  \item[\termblanklong{Interpretation of a two-sample CI}]
        In repeated random sampling with the same sample sizes, approximately \blank{2cm}\%
        of CIs constructed this way will capture \blank{5cm}.
  \item[\termblanklong{Plausible values}]
        Values \blank{3cm} the CI are consistent with the data; values \blank{3cm} the
        CI are not.
  \item[\termblanklong{Width of a CI}]
        $\text{width} = 2 \times t^* \cdot SE$;\quad $SE = \sqrt{\dfrac{s_1^2}{n_1} + \dfrac{s_2^2}{n_2}}$.
        Width \blank{3cm} as sample sizes increase.
\end{description}
\end{vocabbox}

\begin{notesbox}{Part 1: Interpreting a Two-Sample $t$-Interval (UNC-4.Z)}

\textbf{The Interpretation Must Include:}
\begin{enumerate}
  \item The \blank{3cm} level
  \item Both \blank{3cm} (not just the samples)
  \item The \blank{4cm} of subtraction ($\mu_1 - \mu_2$)
  \item The lower and upper bounds with \blank{2cm}
\end{enumerate}

\medskip
\textbf{Model Sentence (2009 AP FRQ 4):}

``Based on these samples, one can be 95 percent confident that the difference in the
population mean response times (northern $-$ southern) is between \blank{2cm} minutes
and \blank{2cm} minutes.''

\medskip
\textbf{Worked Example.} A 99\% CI for $\mu_{\text{online}} - \mu_{\text{in-person}}$ exam
score is $(2.1, 11.3)$ points.

\medskip
Interpretation: ``We are \blank{2cm}\% confident that the true difference in \blank{5cm}
(\blank{2cm} $-$ \blank{2cm}) is between \blank{1.5cm} and \blank{1.5cm} points.''

\end{notesbox}

\begin{notesbox}{Part 2: Justifying a Claim (UNC-4.AA)}

\renewcommand{\arraystretch}{1.7}
\begin{tabularx}{\linewidth}{@{} p{0.35\linewidth} X @{}}
  \textbf{Where is 0?} & \textbf{Conclusion} \\ \hline
  0 is \emph{inside} the interval & \writeline \\
  0 is \emph{above} the interval (entire interval $> 0$) & \writeline \\
  0 is \emph{below} the interval (entire interval $< 0$) & \writeline \\
\end{tabularx}

\medskip
\textbf{Example.} The 99\% CI above is $(2.1, 11.3)$. Is 0 inside or outside?

\quad 0 is \blank{3cm} the interval.

\quad Conclusion: There \blank{3cm} convincing evidence that online students score
higher than in-person students on average.

\end{notesbox}

\begin{notesbox}{Part 3: Effect of Sample Size on Width (UNC-4.AB)}

$SE = \sqrt{\dfrac{s_1^2}{n_1} + \dfrac{s_2^2}{n_2}}$. \quad As $n_1$ and $n_2$ \blank{3cm},
$SE$ \blank{2cm}, so the interval becomes \blank{3cm}.

\medskip
\textbf{Example.} Suppose $s_1 = s_2 = 10$ and we compare two sample-size scenarios:

\renewcommand{\arraystretch}{1.8}
\begin{tabularx}{\linewidth}{@{} p{0.25\linewidth} X X X @{}}
  \textbf{Scenario} & $n_1 = n_2$ & $SE$ & Approx.\ margin of error \\ \hline
  A & 20 & $\sqrt{\dfrac{100}{20}+\dfrac{100}{20}} = \blank{2cm}$ & $\approx$ \blank{2cm} \\
  B & 80 & $\sqrt{\dfrac{100}{80}+\dfrac{100}{80}} = \blank{2cm}$ & $\approx$ \blank{2cm} \\
\end{tabularx}

\medskip
Quadrupling each sample size \blank{3cm} the margin of error by a factor of \blank{1.5cm}.

\end{notesbox}

\begin{practicebox}
A 95\% CI for $\mu_{\text{brand A}} - \mu_{\text{brand B}}$ battery life is $(-3.2,\; 0.4)$ hours.

\begin{itemize}
  \item Interpret the interval: \writelines{2}
  \item Does the interval provide convincing evidence that Brand A lasts longer? Explain.
        \writeline
  \item If both sample sizes were quadrupled (all else equal), would the new CI be
        narrower, wider, or the same? Explain. \writeline
\end{itemize}
\end{practicebox}

\end{document}
